% Provenance (never printed): sources/book2/BRIEFS.md, g10-10. The tax wedge, statutory and economic
% incidence, shares set by slopes and elasticities, deadweight loss growing with the square of the tax,
% revenue against the rate, a subsidy. Real numbers: figdata/grade-10/10-tax-incidence/ (Eurostat HICP,
% Germany, monthly rates). Ledger: German VAT cut 2020, Fuest, Neumeier and Stohlker (2021), Chetty,
% Looney and Kroft (2009). Brenland numbers invented.
\chapter{Tax Incidence and Deadweight Loss}\label{ch:g10:tax-incidence}

On 3 June 2020 the German government announced that value added tax would be cut for six months, from 1 July to 31 December. The
standard rate fell from 19 to 16\,\%, and the reduced rate, which covers most food, from 7 to 5\,\%. The tax is collected by sellers:
a shop adds it to its price and hands it to the state. If shops passed the whole cut on, a price that included 19\,\% of tax would fall
by $3/119$, about 2.5\,\%.
% ledger: g10:de-vat-2020, g10:fns2021 (announcement 3 June 2020)

The opening chart shows the result in Germany's consumer prices. In June 2020 they were 0.8\,\% higher than a year earlier; in July, 0.0\,\%.
In January 2021, when the old rates came back, the twelve-month change jumped from $-0.7$ to 1.6\,\%. Prices did fall, and rise again. But
did they fall by the whole tax cut? And who gains when a tax is cut, or loses when it rises: the buyers, or the sellers? This chapter answers
both questions for a competitive market, and measures what a tax costs beyond the money it raises.
% figdata: vat_steps (cut: all 0.8 -> 0.0; end: all -0.7 -> 1.6)

\begin{openingchart}{Germany: change in consumer prices over twelve months, all items and food, monthly, 2019--2021, in percent. The shaded months
are those of the temporary VAT cut, July to December 2020.}{Eurostat, harmonised index of consumer prices, monthly annual rates of change}
\begin{tikzpicture}
\begin{axis}[omtime, width=0.86\linewidth, height=6cm, xmin=2019, xmax=2022, ymin=-2, ymax=7,
  xtick={2019,2020,2021,2022}, ylabel={percent}, legend pos=north west, legend style={font=\scriptsize}]
\fill[gray!15] (axis cs:2020.5,-2) rectangle (axis cs:2021,7);
\addplot[omDef, thick, mark=*, mark size=0.9pt] table[x=t, y=all] {figdata/grade-10/10-tax-incidence/out/vat_de.csv};
\addlegendentry{all items}
\addplot[omThm, thick, mark=square*, mark size=0.9pt] table[x=t, y=food] {figdata/grade-10/10-tax-incidence/out/vat_de.csv};
\addlegendentry{food}
\draw[gray] (axis cs:2019,0) -- (axis cs:2022,0);
\end{axis}
\end{tikzpicture}
\end{openingchart}

\begin{recall}
From the grades 6--9 volume: a \emph{tax per unit} is a fixed amount of money for each unit sold, and a \emph{subsidy} is a payment per unit. From
\cref{ch:g10:elasticities}: point elasticities of demand and supply. From \cref{ch:g10:surplus-and-efficiency} and \cref{ch:g10:what-prices-do}:
consumer, producer and \omterm{def:g10:surplus-and-efficiency:total-surplus}{total surplus}, and the \omterm{def:g10:what-prices-do:dwl}{deadweight loss}.
\end{recall}

\section{The wedge}\label{sec:g10:tax-incidence:wedge}

A tax of $t$ dalens per unit separates two prices: the price $p_b$ that buyers pay and the price $p_s$ that sellers keep.

\begin{definition}[Tax wedge]\label{def:g10:tax-incidence:wedge}
The \emph{tax wedge}\index{tax wedge} is the gap between the price paid by buyers and the price received by sellers: $p_b - p_s = t$ for a tax of $t$
per unit.
\end{definition}

\begin{definition}[Ad valorem tax]\label{def:g10:tax-incidence:ad-valorem}
An \emph{ad valorem tax}\index{ad valorem tax} is a tax set as a percentage of the price, such as value added tax: with a rate $\tau$ on the seller's
price, $p_b = (1 + \tau)\,p_s$.
\end{definition}

Take the market of \cref{ch:g10:surplus-and-efficiency}: demand $Q = 120 - 4p$, supply $Q = 2p - 24$, an equilibrium of 24 units at 24 dalens. A tax of 6
dalens per unit is collected from the sellers. A seller who must hand over 6 dalens per unit offers the quantity it offered before at the price it now
keeps, $p_s = p_b - 6$: supply becomes $Q = 2(p_b - 6) - 24 = 2p_b - 36$. Then $120 - 4p_b = 2p_b - 36$ gives $p_b = 26$, $p_s = 20$ and $Q = 16$. The tax
raised the buyers' price by 2 and lowered the sellers' price by 4.

\section{Who hands over the tax does not matter}\label{sec:g10:tax-incidence:statutory}

\begin{definition}[Statutory and economic incidence]\label{def:g10:tax-incidence:statutory}
The \emph{statutory incidence}\index{statutory incidence} of a tax says who must, by law, hand the money to the state. Its \emph{economic
incidence}\index{economic incidence} says how the burden is shared: by how much the buyers' price rises and the sellers' price falls.
\end{definition}
\phantomsection\label{def:g10:tax-incidence:economic}

\begin{proposition}[Statutory incidence does not change economic incidence]\label{prop:g10:tax-incidence:same}
In a competitive market, a tax of $t$ per unit collected from sellers and the same tax collected from buyers lead to the same buyers' price, the same sellers'
price and the same quantity.
\end{proposition}
\begin{proof}
Whoever hands over the tax, buyers choose their quantity by the price they pay in all, $p_b$, and sellers by the price they keep, $p_s$, with $p_b = p_s + t$.
The equilibrium is the pair of prices for which $Q^d(p_b) = Q^s(p_s)$ and $p_b - p_s = t$. These two conditions do not mention who hands over the tax, so
they have the same solution in both cases.
\end{proof}

\begin{omfigure}
\begin{tikzpicture}[econ, x=0.12cm, y=0.15cm]
\foreach \side/\sh in {sellers/0, buyers/7.4} {
\begin{scope}[xshift=\sh cm]
\draw[econ axes] (0,33) node[left, font=\small] {$p$} |- (50,0) node[below left=12pt and 0pt, font=\small] {$Q$};
\draw[demand] (0,30) -- (48,18);
\draw[supply] (0,12) -- (40,32);
\draw[guide] (0,26) -- (16,26) -- (16,0); \draw[guide] (0,20) -- (16,20);
\node[eqpoint] at (24,24) {};
\foreach \y in {20,24,26} \draw (0,\y) -- (-0.6,\y) node[left, font=\footnotesize] {\y};
\draw (16,0) -- (16,-0.5) node[below, font=\footnotesize] {16};
\draw (24,0) -- (24,-0.5) node[below, font=\footnotesize] {24};
\node[font=\scriptsize, anchor=north] at (25,-3.4) {tax collected from \side};
\end{scope}}
\begin{scope}
\draw[supply, shifted] (0,18) -- (30,33);
\node[font=\scriptsize, text=omThm!80!black, anchor=west] at (30,32) {supply $+\,6$};
\end{scope}
\begin{scope}[xshift=7.4cm]
\draw[demand, shifted] (0,24) -- (48,12);
\node[font=\scriptsize, text=omDef!80!black, anchor=north east] at (48,11.5) {demand $-\,6$};
\end{scope}
\end{tikzpicture}

{\small A tax of 6 dalens per unit. Left: collected from sellers, it moves the supply curve up by 6. Right: collected from buyers, it moves the demand curve down
by 6. In both cases 16 units are traded, buyers pay 26 and sellers keep 20.}
\end{omfigure}

\section{The less elastic side bears more}\label{sec:g10:tax-incidence:shares}

\begin{definition}[Pass-through rate]\label{def:g10:tax-incidence:pass-through}
The \emph{pass-through rate}\index{pass-through rate} of a tax is the rise in the price paid by buyers divided by the tax: the share of the tax that buyers bear.
\end{definition}

\begin{proposition}[Shares of the burden]\label{prop:g10:tax-incidence:shares}
With demand $Q = a - b\,p$ and supply $Q = c + d\,p$ ($b, d > 0$), a tax of $t$ per unit raises the buyers' price by $\dfrac{d}{b + d}\,t$, lowers the sellers'
price by $\dfrac{b}{b + d}\,t$, and lowers the quantity by $\dfrac{b\,d}{b + d}\,t$. The \omterm{def:g10:tax-incidence:pass-through}{pass-through rate} is $d/(b + d)$.
\end{proposition}
\begin{proof}
With $p_s = p_b - t$: $a - bp_b = c + d(p_b - t)$, so $p_b = \dfrac{a - c + dt}{b + d} = p^* + \dfrac{d}{b + d}\,t$, by \cref{prop:g10:market-equilibrium:solution}.
Then $p_s = p_b - t = p^* - \dfrac{b}{b + d}\,t$, and $Q = a - bp_b = Q^* - \dfrac{bd}{b + d}\,t$.
\end{proof}

In the example, $b = 4$ and $d = 2$: buyers bear $2/6$ of the tax, 2 dalens, and sellers $4/6$, 4 dalens, as found above. Demand is the side that reacts more
to price here (4 units per dalen against 2), and buyers escape more of the tax by buying less.

The same shares can be written with elasticities. At the equilibrium, $|\varepsilon_d| = b\,p^*/Q^*$ and $\varepsilon_s = d\,p^*/Q^*$, so
\[
  \text{pass-through rate} = \frac{d}{b + d} = \frac{\varepsilon_s}{\varepsilon_s + |\varepsilon_d|}.
\]
The side whose elasticity is smaller bears the larger share. If supply is \omterm{def:g10:elasticities:perfectly-elastic}{perfectly inelastic}, as for land, whose quantity no tax can change, $\varepsilon_s = 0$ and
owners bear the whole tax. If demand is \omterm{def:g10:elasticities:perfectly-elastic}{perfectly inelastic}, buyers do.

\section{Deadweight loss and revenue}\label{sec:g10:tax-incidence:dwl}

The state collects $t$ on each unit still traded, $R = t\,Q_t$. Buyers and sellers lose more than that: the units that are no longer traded carried surplus too.

\begin{proposition}[The deadweight loss grows with the square of the tax]\label{prop:g10:tax-incidence:square}
In the linear market of \cref{prop:g10:tax-incidence:shares}, a tax of $t$ per unit causes a \omterm{def:g10:what-prices-do:dwl}{deadweight loss}
\[
  \mathrm{DWL} = \tfrac12\, t\,(Q^* - Q_t) = \tfrac12\,\frac{b\,d}{b + d}\,t^2 .
\]
Doubling the tax multiplies the \omterm{def:g10:what-prices-do:dwl}{deadweight loss} by four.
\end{proposition}
\begin{proof}
The units from $Q_t$ to $Q^*$ are no longer traded. Their lost surplus is the triangle between the demand and supply lines from $Q_t$ to $Q^*$: its vertical side at
$Q_t$ is the wedge $p_b - p_s = t$ and its width is $Q^* - Q_t = \frac{bd}{b+d}\,t$ (\cref{prop:g10:tax-incidence:shares}), so its area is $\tfrac12\,t \cdot
\frac{bd}{b+d}\,t$. Replacing $t$ by $2t$ multiplies $t^2$ by 4.
\end{proof}

In the example, $bd/(b + d) = 8/6 = 4/3$. A tax of 3 dalens cuts the quantity by 4 and costs $\tfrac12 \times 3 \times 4 = 6$ dalens of \omterm{def:g10:what-prices-do:dwl}{deadweight loss}; a tax of 6 cuts
it by 8 and costs $\tfrac12 \times 6 \times 8 = 24$, four times as much, while revenue goes from $3 \times 20 = 60$ to $6 \times 16 = 96$, less than twice as much.

\begin{definition}[Laffer curve]\label{def:g10:tax-incidence:laffer}
The \emph{Laffer curve}\index{Laffer curve} is the name given to the graph of a tax's revenue against its rate.
\end{definition}

For the linear market, $R(t) = t\,\big(Q^* - \tfrac{bd}{b+d}\,t\big)$, a parabola open downward, zero at $t = 0$ and at the tax that stops all trade. In the example,
$R(t) = t(24 - \tfrac43 t)$ is largest at $t = 9$, where 12 units are traded and the state collects 108 dalens; the \omterm{def:g10:what-prices-do:dwl}{deadweight loss} there is $\tfrac12 \times 9 \times 12
= 54$. Beyond 9 dalens, a higher tax raises less money.

\begin{omfigure}
\begin{tikzpicture}[econ, x=0.12cm, y=0.15cm]
\begin{scope}
\draw[econ axes] (0,33) node[left, font=\small] {$p$} |- (50,0) node[below left=12pt and 0pt, font=\small] {$Q$};
\fill[omProp!25] (0,20) rectangle (16,26);
\fill[gray!35] (16,26) -- (24,24) -- (16,20) -- cycle;
\draw[demand] (0,30) -- (48,18);
\draw[supply] (0,12) -- (40,32);
\draw[guide] (16,20) -- (16,0);
\node[font=\scriptsize] at (8,23) {revenue 96};
\node[font=\scriptsize, anchor=west] at (25,21) {DWL 24};
\draw[gray] (19,23) -- (25.2,21.2);
\foreach \y in {20,26} \draw (0,\y) -- (-0.6,\y) node[left, font=\footnotesize] {\y};
\draw (16,0) -- (16,-0.5) node[below, font=\footnotesize] {16};
\end{scope}
\begin{scope}[xshift=7.4cm, x=0.24cm, y=0.034cm]
\draw[econ axes] (0,125) node[left, font=\small] {$R$} |- (20,0) node[below left=12pt and 0pt, font=\small] {tax $t$};
\draw[curve, domain=0:18, samples=50] plot (\x, {\x*(24 - 4*\x/3)});
\draw[guide] (0,108) -- (9,108) -- (9,0);
\draw[guide] (0,96) -- (6,96) -- (6,0);
\node[eqpoint] at (9,108) {};
\foreach \y in {96,108} \draw (0,\y) -- (-0.25,\y) node[left, font=\footnotesize] {\y};
\foreach \x in {6,9,18} \draw (\x,0) -- (\x,-3) node[below, font=\footnotesize] {\x};
\end{scope}
\end{tikzpicture}

{\small Left: a tax of 6 dalens raises 96 (the rectangle: 6 dalens on each of 16 units) and destroys 24 of surplus (the triangle). Right: revenue against the tax per unit. It
peaks at 108, for a tax of 9, and falls to zero at 18, where nothing is traded.}
\end{omfigure}

\section{A subsidy}\label{sec:g10:tax-incidence:subsidy}

A subsidy of $s$ per unit is a negative tax: sellers keep $p_s = p_b + s$. In the example, a subsidy of 6 gives $120 - 4p_b = 2(p_b + 6) - 24$, so $p_b = 22$, $p_s = 28$ and $Q =
32$. The state pays $6 \times 32 = 192$ dalens. Buyers gain 2 dalens per unit and sellers 4, the same shares as for the tax. But the 8 extra units cost sellers more to produce
than buyers value them: the gains of buyers and sellers add up to $192 - 24 = 168$, and 24 dalens of the state's money buy nothing.

\begin{model}{Tax incidence in a competitive market}
\assumes The competitive market of \cref{ch:g10:market-equilibrium}; buyers and sellers react to the prices they pay and keep, including the tax.
\predicts A tax raises the buyers' price by less than the tax when both curves slope, by a share $\varepsilon_s/(\varepsilon_s + |\varepsilon_d|)$, whoever hands it over; a tax cut
  lowers it by the same share; the \omterm{def:g10:what-prices-do:dwl}{deadweight loss} grows with the square of the tax.
\tested The German VAT cut of 2020 (record below): prices fell, by about 70\,\% of the cut in supermarkets, less than fully, as the model predicts when supply slopes.
\breaks When sellers are few enough to set their prices (grade 11); when buyers do not notice a tax that is not shown in the price: a grocery store that showed prices with tax
  included sold 8\,\% less of the goods concerned (Chetty, Looney and Kroft 2009), so the same tax worked differently depending on how it was shown.
\end{model}
% ledger: g10:fns2021, g10:clk2009

\begin{record}[The German VAT cut of 2020]
% ledger: g10:de-vat-2020, g10:fns2021; figdata: vat_steps
From 1 July to 31 December 2020, Germany cut its VAT rates from 19 to 16\,\% and from 7 to 5\,\%. A study of about \num{130000} supermarket products, whose prices were collected
online every day and compared with the same products in Austria, where the tax did not change, found that prices fell by roughly 1.3\,\%: about 70\,\% of the cut reached buyers
(Fuest, Neumeier and Stöhlker 2021). When the rates went back up in January 2021, the price effect was only about half as large. The model predicts the first finding: a cut is
shared. It does not predict the second, since in the model a rise and a fall of the same size move prices by the same amount. The chart shows a further trap for readers of
twelve-month changes: in July 2021, a year after the cut, the change jumped again, from 2.1 to 3.1\,\%, because prices were being compared with the low prices of July 2020.
\end{record}
% figdata: vat_de (2021-06 all 2.1, 2021-07 all 3.1)

\begin{measure}[How pass-through is measured]
To measure how much of a tax change reaches buyers, compare the prices of the goods it touched, just before and just after it, with the prices of similar goods that it did not
touch over the same days: the same products in a neighbouring country, or goods taxed at a rate that did not change. The difference between the two changes, divided by the change in
the tax, is the \omterm{def:g10:tax-incidence:pass-through}{pass-through rate}. Daily prices matter: over months, costs, demand and other taxes move too. A twelve-month rate of change, like those of the opening chart, mixes the
tax with everything else that happened in the year.
\end{measure}

\section{Exercises}

\begin{exercise}[$\star$]\label{exo:g10:tax-incidence:1}
A tax of 2 dalens per unit is collected from sellers. After the tax, buyers pay 15 dalens. What do sellers keep? What is the wedge?
\end{exercise}

\begin{exercise}[$\star$]\label{exo:g10:tax-incidence:2}
Demand is $Q = 100 - 2p$ and supply $Q = 3p - 50$. A tax of 5 per unit is collected from sellers. Find the buyers' price, the sellers' price and the quantity.
\end{exercise}

\begin{exercise}[$\star$]\label{exo:g10:tax-incidence:3}
Read the opening chart. Give the all-items twelve-month change in June and in July 2020, and in December 2020 and January 2021. Which months saw the largest jumps?
\end{exercise}

\begin{exercise}[$\star$]\label{exo:g10:tax-incidence:4}
In \cref{exo:g10:tax-incidence:2}, what share of the tax do buyers bear? Check with \cref{prop:g10:tax-incidence:shares}.
\end{exercise}

\begin{exercise}[$\star$]\label{exo:g10:tax-incidence:5}
In \cref{exo:g10:tax-incidence:2}, how much revenue does the tax raise?
\end{exercise}

\begin{exercise}[$\star\star$]\label{exo:g10:tax-incidence:6}
Compute the \omterm{def:g10:what-prices-do:dwl}{deadweight loss} of the tax in \cref{exo:g10:tax-incidence:2}.
\end{exercise}

\begin{exercise}[$\star\star$]\label{exo:g10:tax-incidence:7}
Demand is $Q = 60 - p$ and supply $Q = p$. A tax of 10 per unit is collected from buyers. Find the two prices, the quantity, the revenue and the \omterm{def:g10:what-prices-do:dwl}{deadweight loss}.
\end{exercise}

\begin{exercise}[$\star\star$]\label{exo:g10:tax-incidence:8}
Demand is $Q = 48 - p$ and supply $Q = 2p - 24$. Find the equilibrium, then the effect of a tax of 6 per unit. Compare the buyers' share with the example of the text, and explain
the difference.
\end{exercise}

\begin{exercise}[$\star\star$]\label{exo:g10:tax-incidence:9}
In the market of the text (demand $Q = 120 - 4p$, supply $Q = 2p - 24$), an \omterm{def:g10:tax-incidence:ad-valorem}{ad valorem tax} of 50\,\% of the seller's price is collected from sellers. Find $p_s$, $p_b$, the quantity and the
revenue.
\end{exercise}

\begin{exercise}[$\star\star$]\label{exo:g10:tax-incidence:10}
In the market of \cref{exo:g10:tax-incidence:7}, a subsidy of 10 per unit is paid instead. Find the two prices and the quantity, the cost to the state and the \omterm{def:g10:what-prices-do:dwl}{deadweight loss}.
\end{exercise}

\begin{exercise}[$\star\star$]\label{exo:g10:tax-incidence:11}
In the market of \cref{exo:g10:tax-incidence:7}, the tax is raised to 20. Compute the new \omterm{def:g10:what-prices-do:dwl}{deadweight loss} and revenue, and compare with a tax of 10.
\end{exercise}

\begin{exercise}[$\star\star\star$]\label{exo:g10:tax-incidence:12}
\Cref{prop:g10:tax-incidence:same} says that it does not matter who hands over the tax. What evidence would contradict it? Use the study cited in the model box.
\end{exercise}

\begin{exercise}[$\star\star\star$]\label{exo:g10:tax-incidence:13}
Show that, in the linear market, the revenue-maximising tax is $t^* = Q^*(b + d)/(2bd)$, and that it cuts the quantity to half of $Q^*$. Apply it to the market of
\cref{exo:g10:tax-incidence:7}: find $t^*$, the revenue and the \omterm{def:g10:what-prices-do:dwl}{deadweight loss}.
\end{exercise}

\begin{exercise}[$\star\star\star$]\label{exo:g10:tax-incidence:14}
A town taxes land at 100 dalens per hectare a year. The quantity of land is fixed. Who bears the tax? What is the \omterm{def:g10:what-prices-do:dwl}{deadweight loss}? Would the answer change for a tax on buildings?
\end{exercise}

\begin{exercise}[$\star\star\star$]\label{exo:g10:tax-incidence:15}
A sales tax is added at the till and not shown on the shelf. Explain, using the record of Chetty, Looney and Kroft (2009), why buyers might react to it less than to the same rise in a
shelf price, and why this breaks the reasoning of \cref{prop:g10:tax-incidence:same}.
\end{exercise}

\section{Problem: Brenland's Fuel Tax}

\begin{problem}[{Weekend problem --- a fuel tax, its incidence, its revenue and its cost}]
\label{pb:g10:tax-incidence:1}
In Brenland, with $p$ the price of fuel in dalens per litre and quantities in millions of litres a year, the demand for fuel is $Q = 930 - 300\,p$ and the supply $Q = 600\,p - 420$.
Revenues and losses are in millions of dalens a year.

\medskip
\noindent\textbf{Part I --- Incidence.}
\begin{enumerate}
\item Find the equilibrium price and quantity without tax.
\item A tax of 0.20 dalens per litre is collected from the fuel sellers. Find the buyers' price, the sellers' price and the quantity.
\item What share of the tax do drivers bear? Check with \cref{prop:g10:tax-incidence:shares}.
\item Compute the point elasticities of demand and supply at the equilibrium without tax, and the share $\varepsilon_s/(\varepsilon_s + |\varepsilon_d|)$.
\item What would change if the tax were collected from drivers at the pump instead?
\end{enumerate}

\medskip
\noindent\textbf{Part II --- Revenue and deadweight loss.}
\begin{enumerate}[resume]
\item Compute the revenue of the tax of 0.20.
\item Compute its \omterm{def:g10:what-prices-do:dwl}{deadweight loss}.
\item The tax is raised to 0.40. Find the two prices, the quantity, the revenue and the \omterm{def:g10:what-prices-do:dwl}{deadweight loss}.
\item By what factor did the \omterm{def:g10:what-prices-do:dwl}{deadweight loss} grow, and by what factor the revenue?
\end{enumerate}

\medskip
\noindent\textbf{Part III --- The revenue curve.}
\begin{enumerate}[resume]
\item Write the revenue $R(t)$ as a function of the tax per litre.
\item Find the tax that brings the most revenue.
\item Find the quantity, the buyers' price and the sellers' price at that tax.
\item Compute the revenue and the \omterm{def:g10:what-prices-do:dwl}{deadweight loss} at that tax.
\item Compute the \omterm{def:g10:what-prices-do:dwl}{deadweight loss} per dalen of revenue at 0.20, at 0.40 and at the revenue-maximising tax.
\end{enumerate}

\medskip
\noindent\textbf{Part IV --- Ten years later.}
\begin{enumerate}[resume]
\item Over ten years, drivers buy more frugal cars, and demand becomes $Q = 1380 - 600\,p$. Check that the equilibrium without tax is unchanged.
\item With this demand, find the quantity, revenue and \omterm{def:g10:what-prices-do:dwl}{deadweight loss} of a tax of 0.40.
\item What share of the tax do drivers bear now?
\item Find the revenue-maximising tax with this demand and the revenue it brings.
\item Explain why a more elastic demand lowers the revenue a fuel tax can raise and raises its \omterm{def:g10:what-prices-do:dwl}{deadweight loss}.
\item State the result of the problem: the revenue-maximising tax per litre and the \omterm{def:g10:what-prices-do:dwl}{deadweight loss} at that rate, with the demand of Part I.
\end{enumerate}
\end{problem}
